\subsection{Уравнение переноса излучения}

Объемный рендеринг строится на дифференциальном уравнении, описывающем поведения света в среде и называемым Уравнением переноса излучения. В наиболее общем виде оно формулируется следующим образом:
%
\begin{equation}
	\label{eq:phys_rte}
	\frac{\mathrm dI_\nu}{\mathrm ds} = - \alpha_\nu I_\nu + \varepsilon_\nu,
\end{equation}
%
где $I_\nu$ - текущая интенсивность излучения на частоте $\nu$, $\alpha_\nu$ - коэффициент поглощения, $\varepsilon_\nu$ - коэффициент излучения.

Численным решением уравнения занимаются различные движки рендеринга, например работа \cite{pbrt}, которая использовалась для написания теории настоящего раздела.

В компьютерной графике принято интенсивность излучения обозначать за $L$, а преодоленное расстояние за $t$. Не ограничивая общности, далее рассматривается конкретная частота $\nu$, которая в формулах будет опускаться.

Частица света, попадая в объем, может встретить частицу среды и:
%
\begin{enumerate}
	\item быть поглощенной ею (absorption);
	\item перенаправленной в другую или с некоторой вероятностью в ту же сторону (scattering);
	\item частица света вообще может не встретиться с объемом и пролететь дальше без изменений.
\end{enumerate}
%
В точке $\mathrm p$ свет, приходящий из направления $-\omega$ и равный $L_i(p, -\omega)$, в силу поглощения выйдет с интенсивностью $L_o(p, \omega)$ не весь. Описать изменение интенсивности света в силу поглощения можно с помощью величины коэффициента поглощения: $\sigma_a$:
%
\begin{equation}
	\label{eq:dl_absorption}
	\mathrm dL(\mathrm p, \omega) = L_o(\mathrm p, \omega) - L_i(\mathrm p, -\omega) = -\sigma_a(\mathrm p, \omega) L_i(\mathrm p, -\omega) \mathrm dt.
\end{equation}
%
Каждая поглощенная частица отдает объему какую-то энергию, и для сохранения равновесия объем может сам начать излучать свет. Таким образом, влияние собственного излучения объема на итоговый свет может быть выражено как:
%
\begin{equation}
	\label{eq:dl_emission}
	\mathrm dL(p, \omega) = \sigma_a(\mathrm p, \omega) L_e(\mathrm p, \omega) \mathrm dt,
\end{equation}
%
где $L_e$ это также собственное свойство объёма, задаваемое зараннее, либо в художественных целях можно задать, что $\mathrm dL(p, \omega) = \mathrm{const}$ (здесь подразумевается вклад только собственного излучения).

Упомянутое равнее рассеяние несет вклад в итоговый свет с двух противоположных сторон. Если поступающее излучение $L_i$ подвергается рассеянию, оно перенаправляется объемом по некоторому распределению $p(\omega, \omega')$. $p$ называется фазовой функцией и задает для каждого входящего света с направления $\omega$ плотность распределения на сфере $S^2$ для выходящих направлений $\omega'$, то есть:
\[
	\int_{S^2} p(\omega, \omega') \mathrm d\omega' = 1.
\]
С одной стороны, если бы объема не было, любой приходящий свет $L_i(\mathrm p, -\omega)$ мог быть нетронутым, но в силу рассеяния, может изменить свое направление, отрицательно повлияв на выходящий свет $L_o(\mathrm p, \omega)$. Это можно описать с помощью коэффициента рассеяния $\sigma_s$:
%
\begin{equation}
	\label{eq:dl_outscattering}
	\mathrm dL(\mathrm p, \omega) = -\sigma_s(\mathrm p, \omega) L_i(\mathrm p, \omega) \mathrm dt
\end{equation}
%
Это же явление может перенаправить весь входящий с других направлений $\omega_i$ свет в рассматриваемую сторону, положительно повлияв на выходящий свет $L_o(\mathrm p, \omega)$:
%
\begin{equation}
	\label{eq:dl_inscattering}
	\mathrm dL(\mathrm p, \omega) = \sigma_s(\mathrm p, \omega) \left( \int_{S^2} p(\mathrm p, \omega_i, \omega) L_i(\mathrm p, \omega_i) \mathrm d\omega_i \right) \mathrm dt
\end{equation}
%
С наблюдениями \ref{eq:dl_absorption}, \ref{eq:dl_emission}, \ref{eq:dl_inscattering}, \ref{eq:dl_outscattering}, уравнение \ref{eq:phys_rte} может быть переписано в следующей форме:
%
\begin{equation}
	\label{eq:graphics_rte}
	\frac{\mathrm dL}{\mathrm dt} = -\sigma_t L + \sigma_s L_s + L_e,
\end{equation}
%
где $\sigma_t = \sigma_a + \sigma_s$ (коэффициент экстинкции, описывающий вероятность встретить частицей света частицу среды, или ещё это означает, что в среднем частице понадобится пролететь расстояние $\sigma_t^{-1}$ чтобы встретить частицу среды), $L_s$ соответствует интегралу из уравнении \ref{eq:dl_inscattering}, а $L_e$ уже содержит в себе $\sigma_a$ из уравнения \ref{eq:dl_emission}.

Подходы к решению уравнения, от быстрых, но неточных решений, до медленных, но сильно приближенных к точному решению, будут описаны в оставшихся разделах про прямой рендеринг.